\documentclass{article}
\usepackage{amsmath, amssymb, amsthm}
\usepackage{hyperref}
\usepackage{geometry}
\geometry{margin=1in}

\title{Erd\H{o}s Problem \#1131}
\date{Last updated: 2026-01-01}
\author{Source: erdosproblems.com}

\begin{document}
\maketitle

\noindent\textbf{Status:} OPEN \quad \textbf{No prize}

\noindent\textbf{Tags:} analysis, polynomials

\medskip
\noindent\textbf{Problem Statement:}

For $x_1,\ldots,x_n\in [-1,1]$ let\[l_k(x)=\frac{\prod_{i\neq k}(x-x_i)}{\prod_{i\neq k}(x_k-x_i)},\]which are such that $l_k(x_k)=1$ and $l_k(x_i)=0$ for $i\neq k$.

What is the minimal value of\[I(x_1,\ldots,x_n)=\int_{-1}^1 \sum_k \lvert l_k(x)\rvert^2\mathrm{d}x?\]In particular, is it true that\[\min I =2-(1+o(1))\frac{1}{n}?\]

\bigskip
\noindent\textbf{Remarks:}

Erd\H{o}s first conjectured this minimum was achieved by taking the $x_i$ to be the roots of the integral of the Legendre polynomial, since Fejer \cite{Fe32} had earlier shown these to be minimisers of\[\max_{x\in [-1,1]}\sum_k \lvert l_k(x)\rvert^2.\]This was disproved by Szabados \cite{Sz66} for every $n>3$. 

Erd\H{o}s, Szabados, Varma, and V\'{e}rtesi \cite{ESVV94} proved that\[2-O\left(\frac{(\log n)^2}{n}\right)\leq \min I\leq 2-\frac{2}{2n-1}\]where the upper bound is witnessed by the roots of the integral of the Legendre polynomial as above.

\bigskip
\noindent\textbf{References:}

[ESVV94] Erd\H{o}s, P. and Szabados, J. and Varma, A. K. and V\'{e}rtesi,
P., On an interpolation theoretical extremal problem. Studia Sci. Math. Hungar. (1994), 55--60.
 
 [Fe32] Fej\'{e}r, Leopold, Bestimmung derjenigen {A}bszissen eines Intervalles, f\"ur
welche die {Q}uadratsumme der {G}rundfunktionen der
{L}agrangeschen Interpolation im Intervalle ein
{M}\"oglichst kleines {M}aximum {B}esitzt. Ann. Scuola Norm. Super. Pisa Cl. Sci. (2) (1932), 263--276.
 
 [Sz66] Szabados, J., On a problem of {P}. {E}rd\H{o}s. Acta Math. Acad. Sci. Hungar. (1966), 155--157.

\bigskip
\noindent\small{Source: \url{https://www.erdosproblems.com/1131}}
\end{document}
